% ============================================================
%  Appendix D: Boolean Lattices and Logical Geometry
% ============================================================

\section*{D.1 Introduction}

Boolean algebra is often introduced as a collection of operators:
$\land, \lor, \neg, \to, \leftrightarrow$.
This approach obscures the underlying geometry.
Logical operators are not isolated symbols — they are points within a
partially ordered space. The Hasse diagrams appearing throughout the monograph
are projections of this deeper structure.

\section*{D.2–D.4 Truth Functions and Encoding}

\begin{definition}{Boolean Function on Two Variables}{bool-func}
  A \emph{Boolean function on two variables} is a map
  $f : \{0,1\}^2 \to \{0,1\}$.
  Since there are four input states and two outputs for each,
  there are $2^4 = 16$ distinct Boolean functions.
\end{definition}

Each function is encoded by its output column over the four input states
$TT, TF, FT, FF$. Thus the space of connectives is $\{0,1\}^4 = Q_4$.

\begin{definition}{Boolean Connective Space}{bool-connective}
  The \emph{Boolean connective space} is $Q_4 = \{0,1\}^4$,
  forming the vertex set of the four-dimensional hypercube.
\end{definition}

\section*{D.5 Partial Order}

\begin{definition}{Truth Inclusion Order}{truth-order}
  For Boolean functions $f, g$, define $f \leq g$ if $f(x) \leq g(x)$
  for every input state $x$.
\end{definition}

\section*{D.7 The Boolean Lattice Theorem}

\begin{theorem}{Boolean Lattice}{bool-lattice}
  The space of two-variable Boolean functions forms a Boolean lattice.
\end{theorem}

\begin{proof}
  $\{0,1\}^4 \cong \mathcal{P}(\{1,2,3,4\})$, with each output pattern
  corresponding uniquely to a subset of input states. The power set forms a
  Boolean lattice under inclusion.
\end{proof}

\begin{figure}[h]
  \centering
  \begin{tikzpicture}[scale=0.85]
    \tikzset{
  boolnode/.style={circle, draw=RSVPdeep, fill=white, minimum size=8mm,
    font=\tiny\ttfamily, line width=0.6pt},
  namednode/.style={circle, draw=RSVPdeep, fill=RSVPlight, minimum size=8mm,
    font=\tiny\bfseries\ttfamily, line width=1.2pt},
  topbot/.style={circle, fill=RSVPdeep, text=white, minimum size=8mm,
    font=\tiny\bfseries\ttfamily, line width=1.2pt},
  edge/.style={RSVPmid!50, thin},
  namedge/.style={RSVPaccent, thick, opacity=0.8},
  lbl/.style={font=\scriptsize\itshape, RSVPmid},
}
% ============================================================
%  diagrams/boolean-hasse-lattice.tex
%  The 16 binary Boolean functions as a Hasse diagram
%  Nodes labeled by 4-bit truth table output
%  Identified named connectives highlighted
% ============================================================

  node distance=1.3cm and 1.1cm,
  boolnode/.style={
    circle, draw=RSVPdeep, fill=white,
    minimum size=8mm, font=\tiny\ttfamily,
    line width=0.6pt
  },
  namednode/.style={
    boolnode, fill=RSVPlight, draw=RSVPdeep, line width=1.2pt,
    font=\tiny\bfseries\ttfamily
  },
  topbot/.style={
    boolnode, fill=RSVPdeep, text=white, line width=1.2pt,
    font=\tiny\bfseries\ttfamily
  },
  edge/.style={RSVPmid!50, thin},
  namedge/.style={RSVPaccent, thick, opacity=0.8},
  lbl/.style={font=\scriptsize\itshape, RSVPmid}
% ── Level 0: Contradiction (bottom) ──────────────────────────
\node[topbot] (n0) at (0,0) {0000};
\node[lbl, below=2pt of n0] {$\bot$};

% ── Level 1: four 1-element functions ─────────────────────────
\node[boolnode] (n1) at (-4.5, 1.6) {0001};   % NOR
\node[boolnode] (n2) at (-1.5, 1.6) {0010};
\node[boolnode] (n4) at ( 1.5, 1.6) {0100};
\node[boolnode] (n8) at ( 4.5, 1.6) {1000};   % AND
\node[lbl, left=2pt of n1] {NOR};
\node[lbl, right=2pt of n8] {AND};

% ── Level 2: six 2-element functions ─────────────────────────
\node[namednode] (n3) at (-5, 3.4) {0011};    % ¬A
\node[boolnode]  (n5) at (-3, 3.4) {0101};
\node[namednode] (n6) at (-1, 3.4) {0110};    % XOR
\node[boolnode]  (n9) at ( 1, 3.4) {1001};    % XNOR
\node[boolnode]  (nA) at ( 3, 3.4) {1010};
\node[namednode] (nC) at ( 5, 3.4) {1100};    % ¬B
\node[lbl, left=2pt of n3] {$\neg A$};
\node[lbl, below right=0pt and 1pt of n6] {XOR};
\node[lbl, right=2pt of nC] {$\neg B$};

% ── Level 3: six 3-element functions ─────────────────────────
\node[namednode] (n7) at (-4.5, 5.2) {0111};  % NAND
\node[boolnode]  (nB) at (-1.5, 5.2) {1011};
\node[namednode] (nD) at ( 1.5, 5.2) {1101};
\node[namednode] (nE) at ( 4.5, 5.2) {1110};  % OR
\node[lbl, left=2pt of n7] {NAND};
\node[lbl, right=2pt of nE] {OR};

% ── Level 4: Tautology (top) ─────────────────────────────────
\node[topbot] (nF) at (0, 6.8) {1111};
\node[lbl, above=2pt of nF] {$\top$};

% ── Edges (level 0 → 1) ───────────────────────────────────────
\foreach \a/\b in {n0/n1, n0/n2, n0/n4, n0/n8}{
  \draw[edge] (\a) -- (\b);
}

% ── Edges (level 1 → 2) ───────────────────────────────────────
\draw[edge] (n1) -- (n3); \draw[edge] (n1) -- (n5);
\draw[edge] (n2) -- (n3); \draw[edge] (n2) -- (n6);
\draw[edge] (n4) -- (n5); \draw[edge] (n4) -- (nC);
\draw[edge] (n8) -- (n9); \draw[edge] (n8) -- (nC);
\draw[edge] (n2) -- (nA); \draw[edge] (n8) -- (nA);
\draw[edge] (n1) -- (n9); \draw[edge] (n4) -- (n6);

% ── Edges (level 2 → 3) ───────────────────────────────────────
\draw[edge] (n3) -- (n7); \draw[edge] (n5) -- (n7);
\draw[edge] (n6) -- (n7); \draw[edge] (n6) -- (nE);
\draw[edge] (n9) -- (nB); \draw[edge] (nA) -- (nB);
\draw[edge] (n3) -- (nB); \draw[edge] (nC) -- (nE);
\draw[edge] (nA) -- (nD); \draw[edge] (n5) -- (nD);
\draw[edge] (n9) -- (nD); \draw[edge] (nC) -- (nB);

% ── Edges (level 3 → 4) ───────────────────────────────────────
\foreach \a in {n7, nB, nD, nE}{
  \draw[edge] (\a) -- (nF);
}

% ── Highlighted path: ⊥ → NOR → ¬A → NAND → ⊤ ──────────────
\draw[namedge] (n0) -- (n1) -- (n3) -- (n7) -- (nF);

% ── Highlighted path: ⊥ → AND → ¬B → OR → ⊤ ─────────────────
\draw[RSVPpurple, thick, opacity=0.7] (n0) -- (n8) -- (nC) -- (nE) -- (nF);

% ── Hypercube level labels ────────────────────────────────────
\node[lbl, right=5.8cm of n0, anchor=west] {Level 0: Contradiction};
\node[lbl, right=5.8cm of n1, anchor=west] {Level 1 (1 true)};
\node[lbl, right=5.8cm of n6, anchor=west] {Level 2 (2 true)};
\node[lbl, right=5.8cm of nE, anchor=west] {Level 3 (3 true)};
\node[lbl, right=5.8cm of nF, anchor=west] {Level 4: Tautology};

% ── Legend ────────────────────────────────────────────────────
\node[draw=RSVPdeep, fill=RSVPlight, rounded corners=3pt,
      inner sep=5pt, font=\scriptsize,
      below right=0.5cm and -1cm of nF, anchor=north] {
  \begin{tabular}{ll}
    \color{RSVPaccent}\textbf{——} & NOR path: $\bot\to\mathrm{NOR}\to\neg A\to\mathrm{NAND}\to\top$\\
    \color{RSVPpurple}\textbf{——} & AND path: $\bot\to\mathrm{AND}\to\neg B\to\mathrm{OR}\to\top$\\
    \color{RSVPdeep}\textbullet\ light & Named connectives
  \end{tabular}
};

  \end{tikzpicture}
  \caption{Hasse diagram of the 16 binary Boolean connectives. Nodes are
    labeled by their 4-bit output pattern. Highlighted paths trace two
    chains from $\bot$ (0000) to $\top$ (1111): the amber path through
    $\mathrm{NOR} \to \neg A \to \mathrm{NAND}$, and the purple path
    through $\mathrm{AND} \to \neg B \to \mathrm{OR}$.}
  \label{fig:hasse-appd}
\end{figure}

\section*{D.8–D.9 Meet, Join, and Complementation}

The lattice possesses meet $f \wedge g$ (intersection of truth sets),
join $f \vee g$ (union of truth sets), and complement $\neg f$ with
$(\neg f)(x) = 1 - f(x)$. Thus $\neg(0110) = 1001$, so XOR and XNOR are
complements.

\section*{D.11 Logic as Geometry}

\begin{definition}{Truth Region}{truth-region}
  For a Boolean function $f$, its \emph{truth region} is
  $R_f = \{x : f(x) = 1\}$.
\end{definition}

Contradiction $\cong$ empty region; tautology $\cong$ full space;
implication $p \to q$ $\cong$ containment $R_p \subseteq R_q$.

\begin{definition}{Logical Distance}{logical-dist}
  The \emph{logical distance} between $f$ and $g$ is
  $d(f,g) = \sum_i |f_i - g_i|$ (Hamming distance).
\end{definition}

\begin{theorem}{Geometry–Logic Correspondence}{geom-logic}
  Every finite Boolean algebra admits a geometric realization as a
  hypercube lattice.
\end{theorem}

\begin{proof}
  A Boolean algebra on $n$ generators is isomorphic to
  $\mathcal{P}(\{1,\ldots,n\}) \cong \{0,1\}^n = Q_n$.
\end{proof}

\begin{namedthm}{The Logic-as-Geometry Principle}
  Logic is not merely symbolic manipulation. Logic is geometry in disguise.
  AND, OR, NAND, XOR, and their kin are not isolated symbols — they are
  vertices of a hypercube ordered by truth content.
\end{namedthm}

\medskip
\textit{Appendix~E} extends this geometric perspective to hypercubes
in arbitrary dimension.
